\documentclass[11pt]{article}

\usepackage[margin=1in]{geometry}
\usepackage{amsmath,amssymb,amsthm}
\usepackage{bm}
\usepackage{booktabs}
\usepackage{float}
\usepackage[colorlinks=true,linkcolor=blue,citecolor=blue,urlcolor=blue]{hyperref}

\newcommand{\avg}[1]{\left\langle #1 \right\rangle}
\newcommand{\Tr}{\operatorname{Tr}}
\renewcommand{\Re}{\operatorname{Re}}
\renewcommand{\Im}{\operatorname{Im}}
\newcommand{\bh}{\bm{h}}
\newcommand{\bF}{\bm{F}}
\newcommand{\bOmega}{\hat{\bm{\Omega}}}
\newcommand{\cG}{\mathcal{G}}
\newcommand{\cF}{\mathfrak{F}}
\newcommand{\ft}[1]{\tilde{#1}}

\newtheorem{proposition}{Proposition}
\theoremstyle{remark}
\newtheorem{remark}{Remark}

\title{The cross-correlation estimator of a CBC background\\
in Stokes form: matrix notation, covariance factorization,\\
and the analytic network-response moment}
\author{astrogwb working note}
\date{\today}

\begin{document}
\maketitle

\begin{abstract}
We rewrite the single-event contribution to the cross-correlation estimator of
Zhong, Iacovelli, Mandic and Berti~\cite{Zhong2026} (their Eqs.~16 and~19) as a
trace over the polarization matrix of the source, expand it in Pauli matrices,
and obtain a contraction of the Stokes four-vector with a four-vector pair
response. For dominant-quadrupole sources whose extrinsic angles are independent
of their intrinsic parameters, the second moment of the estimator factorizes
into a Hadamard product of an intrinsic power moment and an angular response
moment. The angular moment reduces to three constants of the inclination and
three sky integrals, each a finite sum of spherical Bessel functions; a product
quadrature evaluates it exactly, which we check against Monte Carlo draws.
\end{abstract}

%------------------------------------------------------------------------------
\section{The estimator}
%------------------------------------------------------------------------------

Following~\cite{Zhong2026}, the cross-correlation estimator of segment $I$ in a
bin of width $\delta f = 1/\tau$ around $f$ is
\begin{equation}
  \hat C_f^{(I)} = \frac{\kappa_f}{\gamma_{12;f}}
  \Re\!\left[\ft d_{1;f}^{(I)*}\, \ft d_{2;f}^{(I)}\right],
  \qquad
  \kappa_f \equiv \frac{20\pi^2 f^3}{3 H_0^2 \tau},
  \label{eq:estimator}
\end{equation}
with $\gamma_{12;f}$ the isotropic overlap reduction function (ORF), normalized
to one for co-located, co-aligned L-shaped detectors. Neglecting overlapping
chirps and signals that span several segments, the segment average is a
compound-Poisson sum over events plus detector noise~\cite[Eq.~14]{Zhong2026},
\begin{equation}
  \bar C_f = \frac{1}{N_{\rm seg}} \sum_{i=1}^{N_{\rm tot}} X_f^{(i)}
  + \bar n_f,
  \qquad
  X_f^{(i)} = \frac{\kappa_f}{\gamma_{12;f}}
  \Re\!\left[\ft h_{1;f}^{(i)*}\, \ft h_{2;f}^{(i)}\right],
  \label{eq:compound}
\end{equation}
with $N_{\rm tot} \sim \mathrm{Poisson}(N_{\rm seg} R \tau)$.

\paragraph{Detector-space matrix form.}
Stack the two data streams as $\ft{\bm d} = (\ft d_1, \ft d_2)^T$. Since
$\ft{\bm d}^\dagger \sigma_1 \ft{\bm d} = 2 \Re[\ft d_1^* \ft d_2]$,
\begin{equation}
  \hat C_f = \frac{\kappa_f}{2\gamma_{12;f}}\, \ft{\bm d}^\dagger \sigma_1 \ft{\bm d}.
  \label{eq:quadratic}
\end{equation}
For a network of $n$ detectors the optimal combination of all pairs $a<b$ with
weights $w_{ab} = \gamma_{ab}/(S_a S_b)$ is the same quadratic form with
$\sigma_1/\gamma_{12}$ replaced by the symmetric, zero-diagonal $n\times n$ matrix
\begin{equation}
  W_{ab} = W_{ba} = \frac{w_{ab}}{\sum_{c<d} w_{cd}\,\gamma_{cd}}, \qquad
  W_{aa} = 0 .
  \label{eq:network-matrix}
\end{equation}

%------------------------------------------------------------------------------
\section{Polarization matrix and Pauli expansion}
%------------------------------------------------------------------------------

\paragraph{The signal.}
Write the two polarizations and the antenna patterns of detector $a$ as
two-vectors, $\bh = (\ft h_+, \ft h_\times)^T$ and
$\bF_a = (F_a^+, F_a^\times)^T$. Eq.~(17) of~\cite{Zhong2026} reads
\begin{equation}
  \ft h_a = e^{-i\alpha_a}\, \bF_a^T \bh, \qquad
  \alpha_a = 2\pi f\, \bOmega\cdot\bm x_a / c ,
\end{equation}
so the signal in the detector vector is $\ft{\bm d} = \cF \bh$ with the
$n \times 2$ response matrix $\cF_{aA} = e^{-i\alpha_a} F_a^A$.

\paragraph{The polarization matrix.}
The source enters only through the Hermitian, positive semidefinite matrix
\begin{equation}
  \rho \equiv \bh\, \bh^\dagger
  = \begin{pmatrix}
      |\ft h_+|^2 & \ft h_+ \ft h_\times^* \\
      \ft h_\times \ft h_+^* & |\ft h_\times|^2
    \end{pmatrix}.
\end{equation}
With the Stokes parameters of~\cite[Eq.~18]{Zhong2026},
$I = |\ft h_+|^2 + |\ft h_\times|^2$, $Q = |\ft h_+|^2 - |\ft h_\times|^2$,
$U = 2\Re[\ft h_+^* \ft h_\times]$, $V = 2\Im[\ft h_+^* \ft h_\times]$, the
off-diagonal entry is $\ft h_+ \ft h_\times^* = \tfrac12 (U - iV)$ and
\begin{equation}
  \boxed{\;
  \rho = \tfrac12 \left( I\,\mathbf 1 + Q\,\sigma_3 + U\,\sigma_1 + V\,\sigma_2 \right)
  \equiv \tfrac12\, S_\mu\, \sigma^\mu ,
  \qquad
  S_\mu = (I, Q, U, V), \quad
  \sigma^\mu = (\mathbf 1, \sigma_3, \sigma_1, \sigma_2) .
  \;}
  \label{eq:rho}
\end{equation}
Here $Q$ pairs with $\sigma_3$ (the diagonal difference), and $U$, $V$ with the
real and imaginary off-diagonal Pauli matrices. Repeated Greek indices are
summed with no metric; $\Tr[\sigma^\mu \sigma^\nu] = 2\delta^{\mu\nu}$.

\paragraph{Eq.~19 as a trace.}
Inserting $\ft{\bm d} = \cF\bh$ in Eq.~\eqref{eq:quadratic},
\begin{equation}
  \Re\!\left[\ft h_1^* \ft h_2\right]
  = \tfrac12\, \bh^\dagger \cF^\dagger \sigma_1 \cF\, \bh
  = \tfrac12 \Tr\!\left[\rho\; \cF^\dagger \sigma_1 \cF\right],
  \qquad
  \cF^\dagger \sigma_1 \cF = M + M^\dagger, \quad
  M = e^{i\phi_f}\, \bF_1 \bF_2^T ,
  \label{eq:trace}
\end{equation}
with $\phi_f = \alpha_1 - \alpha_2 = 2\pi f\, \bOmega\cdot(\bm x_1 - \bm x_2)/c$, the
phase of~\cite{Zhong2026}. Expanding the Hermitian $2\times2$ matrix
$\tfrac12 \cF^\dagger\sigma_1\cF$ in the same basis,
$\tfrac12\cF^\dagger\sigma_1\cF = \tfrac12 \cG_\mu \sigma^\mu$, and using
$\Tr[\sigma^\mu\sigma^\nu] = 2\delta^{\mu\nu}$:

\begin{proposition}[Four-vector form of Eq.~19]
\begin{equation}
  \boxed{\;
  \Re\!\left[\ft h_1^* \ft h_2\right] = \tfrac12\, S_\mu\, \cG^\mu ,
  \qquad
  \cG^\mu = \tfrac12 \Tr\!\left[\sigma^\mu\, \cF^\dagger \sigma_1 \cF\right]
          = \Re\!\left[ e^{i\phi_f}\, \bF_2^T \sigma^\mu \bF_1 \right] .
  \;}
  \label{eq:four-vector}
\end{equation}
Componentwise, with
$A_I = \bF_1^T \bF_2$, $A_Q = \bF_1^T\sigma_3\bF_2$, $A_U = \bF_1^T\sigma_1\bF_2$
and $A_V = F_1^+F_2^\times - F_1^\times F_2^+$ (so that
$\bF_2^T\sigma_2\bF_1 = i A_V$),
\begin{equation}
  \cG^\mu = \left( \cos\phi_f\, A_I,\; \cos\phi_f\, A_Q,\;
                   \cos\phi_f\, A_U,\; -\sin\phi_f\, A_V \right),
\end{equation}
and the coefficients of~\cite[Eq.~19]{Zhong2026} are
$G_\mu = \tfrac12 \cG^\mu$.
\end{proposition}

The sine reaches only $V$ because $\sigma_2$ is the only antisymmetric Pauli
matrix: the real, symmetric pattern product $\bF_1\bF_2^T + \bF_2\bF_1^T$
cannot couple to it, while the antisymmetric part, which multiplies
$i\sin\phi_f$, couples only to it. A co-located pair ($\phi_f = 0$) is blind to
circular polarization. For a network, $\sigma_1/\gamma_{12}$ becomes $W$ of
Eq.~\eqref{eq:network-matrix} and
$\cG^\mu_{\rm net} = \tfrac12\Tr[\sigma^\mu\cF^\dagger W \cF]
= \sum_{a<b} W_{ab}\, \cG^\mu_{ab}$.

\begin{remark}[Sign conventions]
$\bOmega$ in~\cite{Zhong2026} is the propagation direction. The code
(\texttt{stokes\_cross\_power} in \texttt{astrogwb.detector}) uses the
direction to the source $\hat{\bm n} = -\bOmega$ and
$\phi_{ab} = 2\pi f\,\hat{\bm n}\cdot(\bm x_b - \bm x_a)/c$, which is the same
phase for $(a,b) = (1,2)$.
\end{remark}

\subsection{Structure of the four-vectors}

\paragraph{One source is fully polarized.}
$\rho = \bh\bh^\dagger$ has rank one, so
\begin{equation}
  \det\rho = \tfrac14\left(I^2 - Q^2 - U^2 - V^2\right) = 0 :
\end{equation}
the Stokes vector of every deterministic source is null with respect to
$\eta = \mathrm{diag}(1,-1,-1,-1)$. A sum of sources is partially polarized
(time-like $S$); only in the limit of many sources does it approach the
unpolarized $(I,0,0,0)$ assumed by the sky-averaged ORF. The geometrical shot
noise and polarization leakage of~\cite{Zhong2026} are the finite-catalog
departure from that limit.

\paragraph{Polarization angle.}
A rotation of the polarization frame by $\psi$ acts on $\bh$ as
$R(\psi) = e^{-2i\psi\sigma_2}$, so $\rho \to R\rho R^\dagger$: $I$ and $V$ are
invariant and $(Q,U)$ rotate by $4\psi$. Equivalently it acts on the patterns.
With the complex pattern $\mathbb F_a \equiv F_a^+ + iF_a^\times$, which
transforms as $\mathbb F_a \to e^{\mp 2i\psi}\mathbb F_a$,
\begin{equation}
  A_I + i A_V = \mathbb F_1^*\, \mathbb F_2, \qquad
  A_Q + i A_U = \mathbb F_1\, \mathbb F_2 ,
  \label{eq:complex-patterns}
\end{equation}
so $A_I$, $A_V$ have spin~0 and $A_Q + iA_U$ has spin~4. In terms of the
helicity vector $\bm\epsilon = \hat{\bm m} + i\hat{\bm n}$ of the polarization
frame, $e^+ + ie^\times = \bm\epsilon\otimes\bm\epsilon$ and
$\mathbb F_a = \bm\epsilon^T D_a \bm\epsilon$ with $D_a$ the detector tensor.

\paragraph{The dominant quadrupole.}
For the $(2,2)$ mode in the LAL convention, $\ft h_+ = -\tfrac12(1+c^2)\ft h_0$
and $\ft h_\times = i c\, \ft h_0$ with $c = \cos\iota$, so
\begin{equation}
  S_\mu = |\ft h_0|^2\, s_\mu(\iota), \qquad
  s_\mu = \left( g,\; \tfrac14(1-c^2)^2,\; 0,\; -c(1+c^2) \right), \qquad
  g = \tfrac14(1+c^2)^2 + c^2 ,
  \label{eq:quadrupole}
\end{equation}
null as it must be: $g^2 - s_Q^2 = c^2(1+c^2)^2 = s_V^2$. All frequency
dependence is in $|\ft h_0|^2$, which depends on masses and distance only;
$\psi$ generates $U$ through the rotation above.

%------------------------------------------------------------------------------
\section{Covariance factorization}
%------------------------------------------------------------------------------

\paragraph{Campbell's theorem.}
The $n$-th cumulant of the compound-Poisson sum~\eqref{eq:compound} is
$\mathcal K_n(\bar X_f) = \avg{N_{\rm tot}}\, N_{\rm seg}^{-n}
\avg{(X_f)^n}_\theta$, Eq.~15 of~\cite{Zhong2026}. The same argument for two bins gives the full
frequency covariance, which~\cite{Zhong2026} does not use:
\begin{equation}
  \mathrm{Cov}\!\left[\bar C_f, \bar C_{f'}\right]
  = \frac{\avg{N_{\rm tot}}}{N_{\rm seg}^2}\, \avg{X_f X_{f'}}_\theta
  + \delta_{ff'}\,\sigma_f^2 ,
  \label{eq:campbell}
\end{equation}
with $\sigma_f^2$ the detector-noise variance.

\paragraph{Factorization.}
Split the source parameters into intrinsic ones $\theta_{\rm int}$ (masses,
redshift) and extrinsic ones $\theta_{\rm ext} = (\bOmega, \psi, \iota)$. With
Eqs.~\eqref{eq:four-vector} and~\eqref{eq:quadrupole},
\begin{equation}
  X_f = \underbrace{\frac{\kappa_f}{5}\,\avg{g}\, |\ft h_{0;f}|^2}_{\textstyle P_f(\theta_{\rm int})}
  \;\times\;
  \underbrace{\frac{s_\mu \cG^\mu_{{\rm net};f}}{\avg{s_\mu \cG^\mu_{{\rm net};f}}}}_{\textstyle r_f(\theta_{\rm ext})} ,
  \qquad
  \avg{s_\mu \cG^\mu_{ab}} = \tfrac25\avg{g}\,\gamma_{ab},
  \label{eq:split}
\end{equation}
where $P_f$ is the orientation-averaged contribution of the source and
$\avg{r_f} = 1$. The mean uses that the $\psi$ average removes $A_Q$ and $A_U$
(spin~4), that the $\iota$ average removes $s_V$ (odd in $c$), and that
$\avg{A_I\cos\phi_f}_{\bOmega} = \tfrac25 \gamma_{ab}$, the isotropic ORF in
the normalization above.

\begin{proposition}[Hadamard factorization]
If $\theta_{\rm ext}$ is independent of $\theta_{\rm int}$ and the waveform is
the dominant quadrupole, then
\begin{equation}
  \boxed{\;
  \avg{X_f X_{f'}}_\theta
  = \avg{P_f P_{f'}}_{\theta_{\rm int}} \;
    \avg{r_f r_{f'}}_{\theta_{\rm ext}} ,
  \qquad
  \mathrm{Cov}_{\rm cat} = \frac{\avg{N_{\rm tot}}}{N_{\rm seg}^2}
  \Big( \avg{P P^T} \circ \avg{r r^T} \Big) .
  \;}
  \label{eq:hadamard}
\end{equation}
\end{proposition}

In the code's convention of integrating over the full observation $T$ rather
than averaging segments, the prefactor is $\mathcal R / T$
(\texttt{ImportanceGaussianLikelihood.catalog\_covariance}).

\paragraph{The three components of~\cite{Zhong2026}.}
Writing $r_f - 1$ with Eq.~\eqref{eq:split},
\begin{equation}
  r_f - 1
  = \underbrace{\frac{g\,\cG^I_f}{\avg{g\,\cG^I_f}} - 1}_{\text{geometrical shot noise}}
  \;+\;
  \underbrace{\frac{s_Q\cG^Q_f + s_V\cG^V_f}{\avg{g\,\cG^I_f}}}_{\text{polarization leakage}} ,
\end{equation}
where the first term also carries the inclination dependence of the intensity,
which~\cite{Zhong2026} keep inside $I^{(i)}$ (their ``mean intensity''). The
remaining scatter of $P_f$ between sources, from distances and masses, is the
Poisson shot noise proper. Eq.~\eqref{eq:hadamard} contains all three.

\paragraph{Projection on the mean signal.}
With $c_f = S_f/\sigma_f^2$ and $S_f = \avg{\bar C_f}$, the amplitude estimator
$\hat B = \bm c^T \bar{\bm C} / \bm c^T \bm S$ has relative catalog variance
\begin{equation}
  s^2 = \frac{\bm c^T\, \mathrm{Cov}_{\rm cat}\, \bm c}{(\bm c^T \bm S)^2}
  = \frac{\avg{N_{\rm tot}}}{N_{\rm seg}^2}\,
    \frac{\avg{\left(\textstyle\sum_f c_f P_f r_f\right)^2}}{(\bm c^T\bm S)^2},
  \label{eq:projected}
\end{equation}
the variance that the \texttt{"amplitude"} likelihood adds to $\rho^{-2}$.

\paragraph{When the factorization fails.}
\begin{itemize}
  \item \emph{Higher modes}: the ratios of the Stokes parameters then depend on
        frequency and masses, so $s_\mu$ no longer factors out of $S_\mu$.
  \item \emph{Earth rotation}: a signal long compared with a sidereal day sees
        time-averaged patterns, and the averaging depends on the
        detector-frame chirp mass $(1+z)\mathcal M$ and on $f$, coupling
        $r_f$ to $\theta_{\rm int}$. The static-detector $\avg{r_f r_{f'}}$ is
        then an upper bound at low frequency.
  \item \emph{Selection on SNR}: removing resolved sources by an SNR threshold
        correlates the surviving orientations with distance. A redshift cut
        $z > z_{\min}$, as here and in~\cite{Zhong2026}, does not.
\end{itemize}

%------------------------------------------------------------------------------
\section{The angular moment}
%------------------------------------------------------------------------------

We need, for every pair of pairs $(ab, cd)$ and bins $(f, f')$,
\begin{equation}
  \avg{r_f r_{f'}}
  = \frac{\sum_{a<b}\sum_{c<d} w_{ab}(f)\, w_{cd}(f')\,
          \avg{\Xi_{ab;f}\,\Xi_{cd;f'}}}
         {\left(\tfrac25\avg{g}\right)^2 D_f D_{f'}},
  \qquad
  \Xi_{ab;f} = s_\mu \cG^\mu_{ab;f}, \quad
  D_f = \sum_{a<b} w_{ab}\gamma_{ab} .
\end{equation}

\paragraph{Polarization angle and inclination.}
Since $s_U = 0$,
$\Xi_{ab} = \cos\phi_{ab}\,(g A_I^{ab} + s_Q A_Q^{ab}) - \sin\phi_{ab}\, s_V A_V^{ab}$.
By Eq.~\eqref{eq:complex-patterns} the $\psi$ average removes every product of a
spin-0 and a spin-4 factor, and
\begin{equation}
  \avg{A_Q^{ab} A_Q^{cd}}_\psi
  = \tfrac12 \Re\!\left[\mathbb F_a \mathbb F_b\, \mathbb F_c^* \mathbb F_d^*\right]
  = \tfrac12\left(A_Q^{ab}A_Q^{cd} + A_U^{ab}A_U^{cd}\right).
\end{equation}
The inclination is uniform in $c \in [-1,1]$ and independent of the sky, so its
moments are constants,
\begin{equation}
  \avg{g} = \tfrac45, \quad
  \avg{g^2} = \tfrac{284}{315}, \quad
  \avg{s_Q^2} = \tfrac{8}{315}, \quad
  \avg{s_V^2} = \tfrac{92}{105}, \quad
  \avg{g\, s_V} = \avg{s_Q s_V} = 0 ,
\end{equation}
and $\avg{g\,s_Q} = \tfrac{4}{63}$ multiplies only spin-mixed terms, which vanish.
(The sky-averaged inclination factor of the code is
$\avg{g^2}/\avg{g}^2 = 355/252$.) Therefore
\begin{equation}
  \boxed{\;
  \avg{\Xi_{ab;f}\,\Xi_{cd;f'}}
  = \avg{g^2}\, \Gamma^{II}_{ab,cd}(f,f')
  + \avg{s_Q^2}\, \Gamma^{QQ}_{ab,cd}(f,f')
  + \avg{s_V^2}\, \Gamma^{VV}_{ab,cd}(f,f') ,
  \;}
  \label{eq:three-terms}
\end{equation}
with three ``two-pair ORFs'', sky averages over $\bOmega$ at any fixed $\psi$:
\begin{align}
  \Gamma^{II}_{ab,cd}(f,f') &= \avg{\cos\phi_{ab;f}\cos\phi_{cd;f'}\;
      A_I^{ab} A_I^{cd}}_{\bOmega}, \\
  \Gamma^{QQ}_{ab,cd}(f,f') &= \tfrac12\avg{\cos\phi_{ab;f}\cos\phi_{cd;f'}\;
      \left(A_Q^{ab}A_Q^{cd} + A_U^{ab}A_U^{cd}\right)}_{\bOmega}, \\
  \Gamma^{VV}_{ab,cd}(f,f') &= \avg{\sin\phi_{ab;f}\sin\phi_{cd;f'}\;
      A_V^{ab} A_V^{cd}}_{\bOmega}.
\end{align}
The weights $\avg{g^2} : \avg{s_V^2} : \avg{s_Q^2} \approx 1 : 0.97 : 0.03$
show that the circular-polarization leakage of separated pairs can be as large as
the intensity term, while the linear leakage is small.

\paragraph{Sky integrals in closed form.}
Each pattern is $F_a^A = D_a^{ij} e^A_{ij}(\bOmega,\psi)$. The
$\psi$-invariant products in $\Gamma$ are built from
$\bm\epsilon_i\bm\epsilon_j^* = P_{ij} - i\varepsilon_{ijk}\Omega_k$, with
$P_{ij} = \delta_{ij} - \Omega_i\Omega_j$ the transverse projector (the sign of
the second term is fixed by the handedness of the frame). Every product of four
patterns is therefore a polynomial in $\bOmega$ of degree at most eight,
contracted with $D_a\otimes D_b\otimes D_c\otimes D_d$. The phases combine as
\begin{equation}
  \cos\phi_{ab;f}\cos\phi_{cd;f'}
  = \tfrac12\Re\!\left[e^{i\bm k_-\cdot\bOmega} + e^{i\bm k_+\cdot\bOmega}\right],
  \qquad
  \bm k_\pm = \frac{2\pi}{c}\left(f\,\Delta\bm x_{ab} \pm f'\,\Delta\bm x_{cd}\right),
\end{equation}
and likewise for $\sin\sin$, so each $\Gamma$ is a sum of moments
$\avg{\Omega_{i_1}\cdots\Omega_{i_n} e^{i\bm k\cdot\bOmega}}$ with $n\le8$. These
follow from the generating function
$\avg{e^{i\bm k\cdot\bOmega}} = j_0(k)$:
\begin{equation}
  \avg{\Omega_{i_1}\cdots\Omega_{i_n}\, e^{i\bm k\cdot\bOmega}}
  = (-i)^n\, \partial_{k_{i_1}}\cdots\partial_{k_{i_n}}\, j_0(k)
  = \sum_{p=0}^{\lfloor n/2\rfloor} i^{\,n} (-1)^{p}\,
    \big\{\delta^{p}\,\hat k^{\,n-2p}\big\}_{i_1\cdots i_n}\,
    \frac{j_{n-p}(k)}{k^{p}},
  \label{eq:moments}
\end{equation}
where $\{\delta^p\hat k^{n-2p}\}$ is the sum over the distinct ways of
distributing the indices over $p$ Kronecker deltas and $n-2p$ factors of
$\hat k$, and we used $(k^{-1}\mathrm d/\mathrm dk)^m j_0 = (-1)^m j_m / k^m$.
Contracting with the detector tensors leaves a finite combination of
$j_m(k_\pm)$, $m\le 8$, with coefficients that are scalar invariants of
$D_a, D_b, D_c, D_d$ and $\hat{\bm k}_\pm$. This is the four-detector
generalization of the closed form of the ORF~\cite{Flanagan1993}, which is the
$n\le4$ case of the same construction, and the ground-based analogue of the
four-pulsar average behind the variance of the Hellings--Downs
correlation~\cite{Allen2023}.

\paragraph{Multipole form.}
The sum~\eqref{eq:moments} collapses further. Write the $\psi$-invariant
integrands as
$\Pi^{II} = A_I^{ab}A_I^{cd}$,
$\Pi^{QQ} = \tfrac12(A_Q^{ab}A_Q^{cd} + A_U^{ab}A_U^{cd})$ and
$\Pi^{VV} = A_V^{ab}A_V^{cd}$. Each is a polynomial of degree at most eight in
$\bOmega$, so its angular content stops at $\ell = 8$, and each is even under
$\bOmega\to-\bOmega$ ($A_V$ and $A_U$ change sign, $A_I$ and $A_Q$ do not), so
only even $\ell$ appear. With the plane-wave expansion
$e^{i\bm k\cdot\bOmega} = \sum_\ell (2\ell+1)\, i^\ell j_\ell(k)\,
P_\ell(\hat{\bm k}\cdot\bOmega)$,
\begin{equation}
  \avg{\Pi\, e^{i\bm k\cdot\bOmega}}_{\bOmega}
  = \sum_{\ell = 0,2,4,6,8} (2\ell+1)\,(-1)^{\ell/2}\, j_\ell(k)\,
    \pi_\ell(\hat{\bm k}),
  \qquad
  \pi_\ell(\hat{\bm k}) = \avg{\Pi(\bOmega)\,
    P_\ell(\hat{\bm k}\cdot\bOmega)}_{\bOmega},
  \label{eq:multipole}
\end{equation}
which is real. Combining the two phases as in Eq.~(26), and
$\sin\phi_{ab;f}\sin\phi_{cd;f'} = \tfrac12\Re[e^{i\bm k_-\cdot\bOmega} -
e^{i\bm k_+\cdot\bOmega}]$ for $\Gamma^{VV}$,
\begin{equation}
  \boxed{\;
  \Gamma^{XX}_{ab,cd}(f,f') = \frac12 \sum_{\pm} \epsilon^{XX}_\pm
    \sum_{\ell = 0,2,4,6,8} (2\ell+1)\,(-1)^{\ell/2}\,
    j_\ell(k_\pm)\, \pi^{XX}_\ell(\hat{\bm k}_\pm)
  \;}
  \label{eq:two-pair-orf}
\end{equation}
with $\epsilon^{II}_\pm = \epsilon^{QQ}_\pm = 1$ and
$\epsilon^{VV}_\pm = \mp 1$.
For $\ell \le 4$ and $\Pi = A_I^{ab}$ alone, the same construction gives the
isotropic ORF,
$\gamma_{ab}(f) = \tfrac52 \sum_{\ell=0,2,4}(2\ell+1)(-1)^{\ell/2}\,
j_\ell(k)\, \pi_\ell[A_I^{ab}](\hat{\bm k})$ with
$\bm k = 2\pi f\,\Delta\bm x_{ab}/c$, which is the closed form
of~\cite{Flanagan1993}: two patterns reach $\ell = 4$, four reach $\ell = 8$.

\paragraph{Forty-five numbers per pair of pairs.}
With $\Pi_{\ell m} = \int \mathrm d\bOmega\, \Pi(\bOmega)\, Y^*_{\ell m}(\bOmega)$
and the addition theorem
$P_\ell(\hat{\bm k}\cdot\bOmega) = \frac{4\pi}{2\ell+1}\sum_m
Y_{\ell m}(\hat{\bm k})\,Y^*_{\ell m}(\bOmega)$,
\begin{equation}
  (2\ell+1)\,\pi_\ell(\hat{\bm k}) = \sum_{m=-\ell}^{\ell}
    \Pi_{\ell m}\, Y_{\ell m}(\hat{\bm k}),
\end{equation}
so that
\begin{equation}
  \Gamma^{XX}_{ab,cd}(f,f') = \frac12\sum_\pm \epsilon^{XX}_\pm
  \sum_{\ell\,\rm even}^{8} (-1)^{\ell/2}\, j_\ell(k_\pm)
  \sum_{m} \Pi^{XX}_{\ell m}\, Y_{\ell m}(\hat{\bm k}_\pm).
  \label{eq:harmonic}
\end{equation}
For each pair of pairs and each of $II$, $QQ$, $VV$, the
$1 + 5 + 9 + 13 + 17 = 45$ coefficients $\Pi_{\ell m}$ (45 real numbers, since
$\Pi$ is real and $\Pi_{\ell,-m} = (-1)^m\Pi^*_{\ell m}$) depend only on the
detector tensors. All frequency dependence then sits in $j_\ell(k_\pm)$ and in
the directions $\hat{\bm k}_\pm(f,f')$, so Eq.~\eqref{eq:harmonic} gives the
full $F\times F$ block explicitly.

\paragraph{The same pair.}
For $cd = ab$ (or parallel baselines) $\hat{\bm k}_\pm = \pm\Delta\hat{\bm x}_{ab}$
for every $(f,f')$, and $\pi_\ell$ is even in $\hat{\bm k}$, so the five
multipoles are constants and
\begin{equation}
  \Gamma^{XX}_{ab,ab}(f,f') = \frac12 \sum_\pm \epsilon^{XX}_\pm
  \sum_{\ell = 0,2,4,6,8} (2\ell+1)\,(-1)^{\ell/2}\,
  \pi^{XX}_\ell\big(\Delta\hat{\bm x}_{ab}\big)\;
  j_\ell\!\left(\frac{2\pi |f\pm f'|\,|\Delta\bm x_{ab}|}{c}\right):
\end{equation}
the response correlation of one pair depends on the two bins only through
$f\pm f'$, and decorrelates on the scale $c/|\Delta\bm x_{ab}|$.

\paragraph{Exact multipoles.}
In a frame with polar axis $\hat{\bm k}$, $\Pi\,P_\ell$ has azimuthal harmonics
$|m| \le 8$ and, after the azimuthal average, is a polynomial of degree at
most 16 in $\mu = \hat{\bm k}\cdot\bOmega$. Nine Gauss--Legendre nodes in $\mu$
and $M \ge 9$ trapezoid nodes in azimuth (we use 17) therefore give every
$\pi_\ell$ exactly, independently of $k$. Evaluate Eq.~\eqref{eq:two-pair-orf}
with $j_\ell$ itself (for example \texttt{scipy.special.spherical\_jn}) rather
than its elementary $\sin k$, $\cos k$ form, whose terms up to $k^{-9}$ cancel
catastrophically as $k \to 0$, that is for nearby bins and short baselines.

\paragraph{Validation.}
On the packaged geometry of S2, R2 (ET-2L, misaligned) and C1 (CE Hanford):
\begin{itemize}
  \item the $\ell\le4$ form reproduces
        \texttt{isotropic\_overlap\_reduction\_function} for S2--C1 to twelve
        digits at 5, 30 and 200\,Hz ($+0.232813658228$, $-0.015303275555$,
        $+0.002328551043$), and the $\ell = 6, 8$ multipoles of $A_I$ alone
        vanish to $10^{-17}$;
  \item Eq.~\eqref{eq:two-pair-orf} on the $9\times17$ grid agrees with a
        $200\times400$ brute-force sky quadrature to 10--11 significant
        digits (Table~\ref{tab:two-pair}).
\end{itemize}

\begin{table}[ht]
  \centering
  \small
  \begin{tabular}{lcccc}
    \toprule
    Pairs & $(f, f')$ [Hz] & $\Gamma^{II}$ & $\Gamma^{QQ}$ & $\Gamma^{VV}$ \\
    \midrule
    S2C1 $\times$ S2C1 & $(10, 10)$
      & $+3.3775025974\times10^{-2}$ & $+3.0064142131\times10^{-2}$
      & $+3.9402350621\times10^{-2}$ \\
    S2C1 $\times$ R2C1 & $(10, 17)$
      & $+8.0941043916\times10^{-3}$ & $+1.1493714444\times10^{-3}$
      & $-5.2373201804\times10^{-3}$ \\
    S2R2 $\times$ R2C1 & $(23, 6)$
      & $+9.3111051890\times10^{-4}$ & $+2.8520297874\times10^{-2}$
      & $+8.3486515793\times10^{-4}$ \\
    S2C1 $\times$ R2C1 & $(150, 90)$
      & $-1.9235379575\times10^{-4}$ & $+3.4883349186\times10^{-4}$
      & $+3.9757883079\times10^{-4}$ \\
    \bottomrule
  \end{tabular}
  \caption{Two-pair ORFs from Eq.~\eqref{eq:two-pair-orf}. The brute-force
  quadrature agrees in every digit shown except the last of
  $\Gamma^{QQ}$ at $(150, 90)$\,Hz ($\ldots187$).}
  \label{tab:two-pair}
\end{table}

For S2R2 $\times$ R2C1, $\Gamma^{QQ}$ exceeds $\Gamma^{II}$ thirtyfold: the
misaligned ET pair barely couples to intensity but does couple to linear
polarization, which the small weight $\avg{s_Q^2} = 8/315$ keeps from
dominating Eq.~\eqref{eq:three-terms}.

\paragraph{Exact quadrature.}
In practice the closed form need not be written out. The integrand of
$\avg{r_f r_{f'}}$ is a trigonometric polynomial in $\psi$ with harmonics
$e^{2in\psi}$, $|n|\le 4$, and a polynomial of degree eight in $c=\cos\iota$;
on the sky it is a degree-eight polynomial times plane waves of wavenumber at
most $k_{\max} = 2\pi f_{\max}(|\Delta\bm x_{ab}| + |\Delta\bm x_{cd}|)/c$.
Hence
\begin{itemize}
  \item $K \ge 5$ trapezoid nodes in $\psi\in[0,\pi)$ and $J \ge 5$
        Gauss--Legendre nodes in $c$ are exact (or both averages are done
        analytically as above);
  \item $N$ Gauss--Legendre nodes in $\cos\theta$ times $M$ trapezoid nodes in
        the azimuth converge faster than exponentially once
        $N \gtrsim (k_{\max} + 8)/2$ and $M \gtrsim k_{\max} + 8$.
\end{itemize}
Weighted nodes replace the $L$ isotropic draws of
\texttt{sample\_network\_response}, whose relative error on $\avg{r^2}$ is
$\sqrt{\avg{r^4}/\avg{r^2}^2 - 1}\,/\sqrt{L}$.

\paragraph{Check.}
Table~\ref{tab:check} compares the quadrature with $2^{20}$ isotropic draws for
the SNR-weighted moment $\avg{R^2}$, $R = \sum_f u_f r_f$ with
$u_f = (S_f/\sigma_f)^2/\sum_{f'}(S_{f'}/\sigma_{f'})^2$, which is the
response factor of Eq.~\eqref{eq:projected} if
$\avg{P_fP_{f'}}\propto S_fS_{f'}$. The quadrature uses
$32\times64$ sky nodes, $K = 8$ and $J = 5$, and agrees to all digits shown with
$96\times192$ sky nodes. The draws share seeds across networks, hence one common
fluctuation.

\begin{table}[H]
  \centering
  \small
  \begin{tabular}{lccc}
    \toprule
    Network & quadrature ($8.2\times10^4$ nodes) & Monte Carlo ($2^{20}$ draws)
            & $\avg{R^4}/\avg{R^2}^2$ \\
    \midrule
    ET-2L (misaligned) + CE & $7.245190$ & $7.219 \pm 0.017$ & $6.85$ \\
    ET-2L (aligned) + CE    & $2.276385$ & $2.266 \pm 0.005$ & $5.75$ \\
    ET triangle             & $1.985314$ & $1.978 \pm 0.004$ & $5.39$ \\
    \bottomrule
  \end{tabular}
  \caption{SNR-weighted response second moment. The quadrature takes under a
  second per network; the draws, about twenty.}
  \label{tab:check}
\end{table}

\begin{thebibliography}{9}
\bibitem{Zhong2026}
H.~Zhong, F.~Iacovelli, V.~Mandic and E.~Berti,
``All you need is not $\Omega_{\rm gw}$: beyond the mean of the
cross-correlation estimator when searching for an astrophysical
gravitational-wave background,''
\href{https://arxiv.org/abs/2608.11119}{arXiv:2608.11119} (2026).

\bibitem{Flanagan1993}
\'E.~\'E.~Flanagan,
``Sensitivity of the Laser Interferometer Gravitational Wave Observatory to a
stochastic background, and its dependence on the detector orientations,''
Phys.\ Rev.\ D \textbf{48}, 2389 (1993).

\bibitem{Allen2023}
B.~Allen,
``Variance of the Hellings--Downs correlation,''
Phys.\ Rev.\ D \textbf{107}, 043018 (2023).
\end{thebibliography}

\end{document}
